\section{Experimental design}\label{sec:design}

\subsection{Hypotheses}

The hypotheses are stated here, before the results, with the condition that refutes each. The thresholds of H1 to H3 are those of the research plan of the method and of the acceptance tests of the engine and predate the experiments. The decision rules of H4 to H6 were set after the first measurements of the weight update and of the normalized-error loss were available, so they are treated as post hoc, and the results say so. Table~\ref{tab:hyp} lists them. H1 is the hypothesis of the method: the spread on the crosstalk device is at
least two times below that of phase-fixed WGS at an efficiency at most ten percent lower. H2 is the control that excludes the other explanation of H1, that more iterations or a different optimizer give the gain without the device model. H3 is the fuller acceptance criterion of the research plan, applied to the square lattice and the large lattice of the studio's benchmark: the factor of at least two at an
efficiency ratio of at least $0.9$. It is stricter than H1, which concerns one lattice behind a crosstalk-only device, and it uses a harsher device, with quantization, a fill factor and flicker; the part of the criterion that concerns camera feedback is not tested here. H4 reads the same benchmark by its own win rule at an equal wall-clock time. H5 asks whether the loss is the source of the gain: it states that the efficiency-level loss and the normalized-error loss of Schroff et al.\ return spreads that agree within the seed ranges. H6 is the rationale of the weight update, which the earlier design of the method gave as its integral action, stated in the direction that the ablation found: the update raises the spread.

\begin{table}[t]
\centering\small
\caption{Hypotheses with the condition that refutes each. The thresholds of H1 to H3 are those of the research plan and of the acceptance criteria of the engine; the rules of H4 to H6 are post hoc.}
\label{tab:hyp}
\begin{tabular}{@{}l p{6.4cm} p{6.0cm}@{}}
\toprule
ID & Statement & Refuted if\\
\midrule
H1 & On the crosstalk device the refinement lowers the spread of a $10\times10$ lattice at least two times against phase-fixed WGS, at an efficiency at most ten percent lower & the factor is below two or the efficiency ratio below $0.9$ on a seed\\
H2 & The gain disappears when the device model is removed from the loop & the control keeps a factor of two or more, or more iterations of the baseline close the gap\\
H3 & On the benchmark's device the refinement meets both thresholds of H1 on the square lattice and on the large lattice & either geometry fails either threshold\\
H4 & At the equal wall-clock time of the benchmark the win rule reads the refinement as better than phase-fixed WGS on every geometry & any geometry reads as a tie or as worse\\
H5 & The median spread of the efficiency-level loss lies within the range of the normalized-error loss on the crosstalk device and on the full device & the median lies outside that range on either device\\
H6 & The weight update raises the spread of the refinement on the acceptance task and on the full device & the spread with the update is not higher on either\\
\bottomrule
\end{tabular}
\end{table}

\subsection{Data and provenance}

All data are generated by the engine from seeds; no measured hologram or camera image enters the paper. The lattice task is a $10\times10$ square lattice of equal traps eight FFT bins apart on an SLM of $\parNpix^2$ pixels padded by \parPad, at \parLambda{} nm, with a lens of \parFocal{} mm, a pitch of \parPitch{} $\mu$m and a numerical aperture of
\parNa{} in air, and a Gaussian input beam; it is scored on the crosstalk device (\nSeedsOne{} seeds) and on the full device (\nSeedsOneFull{} seeds). The benchmark geometries are those of the studio's benchmark file: a square lattice (G1), a ring of 36 traps (G2), 100 random traps (G3), a quasi-periodic set of 119 traps (G4) and two large sets, a $32\times32$ lattice (G6a) and a disordered set of
819 traps (G6b), on the device of case C08. The benchmark's three-dimensional geometry is outside the scope of the method and is not used.The benchmark's cells at 100 iterations hold \benchOneDevSeeds{} seeds for the refinement and \benchOneDevSeedsMseven{} for phase-fixed WGS, so they are read as ranges and never as wins; its rows at the time budget of \benchOneTimeBudget{} s hold at least \benchOneWinMinSeeds{} seeds on each side. The smaller grid of the studio's case C08 is an $8\times8$ array on a $\bakeCeightPix^2$-pixel SLM; its baked scores come from
the committed manifest and the committed ranges over \bakeCeightSeedsN{} seeds. The comparison of the weight update and of the normalized-error loss on the four lattices G1 to G4 runs each cell as the benchmark's own record through the engine's runner, scored on the C08 device, for \geoSeeds{} seeds at \geoIters{} outer iterations.

\subsection{Methods compared}

The published methods are run by the product on the same data, not quoted. The baselines are plain GS, weighted GS and phase-fixed WGS, the last at 30 and at 100 iterations, scored on the device; the reference of the comparison is the better of the two phase-fixed runs on the spread and on the efficiency (Eq.~\ref{eq:factor}), because more iterations cannot remove a crosstalk that the model does not contain (Proposition~\ref{prop:floor}).
The ladder methods of the studio are the implementations of the cited algorithms and are checked against the independent library slmsuite~\cite{slmsuite} in the engine's tests. The closest published method, the conjugate-gradient minimization through a crosstalk model of Schroff et al.~\cite{schroff}, is run as the refinement with its loss option set to Eq.~\ref{eq:nel}: the same warm start, device model, conjugate-gradient routine, outer-iteration structure and safeguard, so the comparison isolates the loss. It is our implementation of the loss on point traps, with a closed-form gradient, and not the authors' code, their automatic differentiation, their continuous targets or their camera feedback. The controls are the refinement with the device removed from the loop, the weight update switched on, and the ablated variants of Section~\ref{sec:ablation}.

\subsection{Metrics and statistics}

The metrics are the efficiency and the spread of Eq.~\ref{eq:metrics}, the factor and the efficiency ratio of Eq.~\ref{eq:factor} and the time per hologram. The factor of a seed is a ratio, so its mean and interval are
taken on the logarithmic scale: for $n$ seeds with factors $F_i$,
\begin{equation}
\bar F=\exp\Big(\frac1n\sum_i\log F_i\Big),\qquad \bar F\,e^{\pm t_{n-1}s/\sqrt n},\quad s^2=\frac1{n-1}\sum_i\big(\log F_i-\overline{\log F}\big)^2,\label{eq:geomean}
\end{equation}
with $t_{n-1}$ the two-sided 95 percent quantile of Student's $t$ distribution. The two losses are compared by the median and the range of the spread over the seeds and by the number of seeds on which each is lower, because with five seeds the intervals are wide and the tests that they imply are of low power; no $p$-value is reported, and a claim rests on the range and on the count of seeds. The medians and ranges are
reported as well, because the distributions are not normal (the factor is bounded below by zero and skewed). The benchmark cells report the median and the interquartile range of the studio's protocol, and the reading of a cell follows the win rule of the studio's technical report: a method is better than another in a cell only if both have at least five seeds, their interquartile ranges do not overlap and their medians differ by at least \benchOneWinMinDiff{} percentage points, and ranges that touch count as overlapping; below five seeds, or with overlapping ranges, the cell is a tie, and its $n$ is stated. Rows at a fixed number of iterations give each method its own work and never produce a win. The win rule is one implementation shared with the report, read through the schema check of the benchmark file; here it compares the refinement with phase-fixed WGS, the reference of this paper, and not with the best published method of the cell, so the readings can differ from the report's table.

\subsection{Protocol and ablations}

All methods are deterministic given a seed through the portable generator of the engine, so a rerun reproduces the numbers up to the floating point of the machine's FFT; the times are wall clock and differ between runs. The seeds are $1,\dots,n$ with $n$ stated per experiment: \nSeedsOne{} for the crosstalk device, \nSeedsOneFull{} for the full device, \schroffXN{} and \schroffFN{} for the normalized-error loss on the two devices, and \ablBN{} and \ablBFN{}
for the ablation on the crosstalk device and on the full device. The budget of the refinement is \parOuter{} outer iterations after \parWarm{} warm iterations. The ablation is planned for the components that the method names as its rationale: the weight update, the device-aware conjugate-gradient step, the safeguard, the device model in the loop and the number of steps. Each variant removes one component and keeps the others, on the same seeds and tasks as the main result, and the variant with the default components is checked against the engine's module. The time per hologram of every method is read from its own run on the machine of the experiments, a CPU without an accelerator, and the equal-time comparison of the benchmark gives each method the same wall-clock budget.
